\section{Installing the actual target host on the cohomological windows}
\label{sec:host-window-installation}

\status{HX1--HX4 close a previously unjoined dependency, rather than
claim a new estimate for Galois cohomology.  Relative to the named
classical inputs already recorded for QP and the least-prime theorem
specified below, the literal implication
\texttt{VWinAt}$(\psi,401)\Rightarrow
\texttt{HostExW}(3/2,(\texttt{GChi}\,\psi).\texttt{req})$
is now supplied.  The host is the one constructed from AH18, on the
same prime pair and dyadic field layer that occur in the window input.
The estimate \texttt{VWinAt} itself is not asserted unconditionally.}

\paragraph{The input and output, with their quantifiers.}
Fix a primitive Dirichlet character $\psi$ of conductor $M$ and
order $o$.  The notation $\mathrm{VVle}(E,t,B)$ means both finiteness
of all the relevant $\ell$-primary strict-Selmer lengths in twists
$1\le j<t$, and the bound
\[
 \sum_{\ell\in S}\log\ell
       \sum_{1\le j<t}\operatorname{length}_{\ell}(E,j)\le B
 \quad\text{for every finite set of primes }S.
 \tag{HX.1}
\]
The lengths and local conditions in (HX.1) are the original ones:
the source quantifies over every algebraic closure and every compatible
local-data system.  It is not enough to check one finite set $S$ or one
substitute coefficient model.  Its elementary monotonicity is
\[
 \mathrm{VVle}(E,t,B),\quad t'\le t,\quad B\le B'
       \quad\Longrightarrow\quad\mathrm{VVle}(E,t',B').
 \tag{HX.2}
\]
This follows by restricting the twists and using the nonnegativity of
each length and of $\log\ell$.

The hypothesis $\mathrm{VWinAt}(\psi,401)$ has the following order.
For every $0<\varepsilon<1/52$ and every dyadic page function $e(x)$
with
\[
 x^{2+8\varepsilon}\le2^{e(x)}\le2x^{2+8\varepsilon}
 \qquad(x\ge1),
 \tag{HX.3}
\]
one chooses $J\ge1$ and $X_0$ before every $X\ge X_0$ and every
admissible pair $q,d$ at base $X^{1/J}$.  One then obtains a window
$[W,W^2]\subseteq[X^{1/J},X]$ on which every original target field
$E$ at page $e(x)$ obeys
\[
 \mathrm{VVle}\bigl(E,\lfloor x^{401}\rfloor,
                         \sqrt{\exp(x^\varepsilon)}\bigr).
 \tag{HX.4}
\]
The precise admissible-pair congruences and bounds are the
\texttt{WinPair} record; QP22 supplies the stronger \texttt{PairGood}
record with those same clauses.

The conclusion $\mathrm{HostExW}(3/2,(\mathrm{GChi}\,\psi).\mathrm{req})$
says that, for each $\varepsilon$ in that range, there are constants
$C_{T_2},C_8,C_5,c_6,C_b,x_0$ and one scale set containing arbitrarily
late intervals $[Z,Z^{3/2}]$, together with one fixed $J\ge1$.
The membership condition on a scale $s$ includes $s.J=J$; also
$s.\varepsilon=\varepsilon$.  At every allowed scale $x=s.x\ge x_0$
there are ranks $R$, an actual host $H$, and its auxiliary prime $q_0$ such
that
\[
 \begin{gathered}
 \mathrm{RankVal}(s,R),\qquad H.C_{T_2}=C_{T_2},\qquad
 \mathrm{HostGood}(s,C_8,C_5,c_6,C_b,H,q_0),\\
 \mathrm{EAtH}(\mathrm{Level.cyc}\,\psi,H,TT,x,\varepsilon,4),\qquad
 \mathrm{VVle}(H.E,R.t,\sqrt{s.b}).
 \end{gathered}
 \tag{HX.5}
\]
The scale set and constants precede all of these scale-by-scale
choices.  The hosts may vary with the scale; no single field is
asserted to serve every dyadic page of the interval.

\paragraph{HX1: the exact small-prime input is a classical application.}
The source's \texttt{SmallOne} is
\[
 \exists a\ge1\ \forall N\ge2\ \exists p\text{ prime}:
                p\equiv1\pmod N,\qquad p\le N^a.
 \tag{HX.6}
\]
Xylouris, \emph{On the least prime in an arithmetic progression and
estimates for the zeros of Dirichlet $L$-functions}, Acta Arithmetica
150 (2011), Theorem~1.1, gives an absolute $C>0$ with
$P(a,q)\le Cq^{5.18}$ for coprime positive $a,q$.
Only existence of an absolute exponent is needed here.  Put
$C_+=\max(C,1)$ and
\[
 a_0=\max\left(1,5.18+\frac{\log C_+}{\log2}\right).
\]
For $N\ge2$ the hypotheses hold with residue $1$ and modulus $N$,
and $C_+\le N^{\log C_+/\log2}$.  Hence the supplied prime satisfies
(HX.6) with $a=a_0$.  If the published estimate is read with a
large-modulus threshold, enlarge $C$ over the finitely many smaller
moduli first; Dirichlet's theorem supplies those finitely many primes.
There is no use of RH, GRH, a uniform prime number theorem with a
moving modulus, or a numerical value of the implicit constant.

\paragraph{HX2: ranks and admission are already supplied on this field.}
Suppose AH18 supplies a frame $\phi$ and $H_0:\mathrm{HostAx}(\phi)$,
with
$x^{2+8\varepsilon}\le|\phi|\le2x^{2+8\varepsilon}$.
SZ13 constructs the entire rank record over this exact frame.
Explicitly it chooses even $m,k$ from the two prescribed real rank
scales, sets $f=2|\phi|$, $r=k+m-1$, and defines $t$ by the exact
ceiling expression (SZ.32).  Its proof gives
$R.\mathrm{toFrame}=\phi$ and $\mathrm{RankVal}(s,R)$; it does not
replace $\phi$ by an abstract group of equal cardinality.

For every permitted set of host choices, every
$\mathrm{FieldRows}(H_0,b)$ and $q,d>4tr+1$, HC4 copies the prescribed
CM type, splitting different ideal, square roots and choices into
\texttt{HostNF}.  Thus it gives the full predicate called
\texttt{admitted} in the source, with exact equality of the underlying
\texttt{HostAx} and of the choices.  AH18 supplies the field facts used
in that copying argument.  This proves the complete rank/admission
input of \texttt{admissibleClassesTwistedChar}, uniformly before the
frame and field are selected.  No arithmetic existence assumption is
being hidden in the word ``admitted.''

\paragraph{HX3: all clauses of the conditional host completion.}
Fix $\varepsilon,C_0$ with $0<\varepsilon<1/52$ and $C_0\ge1$.
Use QR25's actual \texttt{ShortNorm} and HX1 in HC5.  The output has
an absolute exponent $C_1$ before $\varepsilon$ and the host.
For every later datum with
\[
 N=Qqd,\qquad d\mid N,\qquad d\ge100,\qquad
 \log N\le3x^{\varepsilon/5}(\log x)^2,
 \tag{HX.7}
\]
it gives the two distinct norm primes $g_0,g_1$, their disjoint
$BB(r,g_i)$ sets, and the three smaller primes $p,p',q_0$ of HC5.
Their norms, integral generators and congruences are (HC.7)--(HC.9),
not merely assertions that the primes split.

Here is the complete input matching for the HC9 host join.
\begin{center}
\begin{tabular}{p{.33\linewidth}p{.58\linewidth}}
Input required on the original host & Supplied construction\\\hline
\texttt{ShortNorm}, \texttt{SmallOne}&QR25 and HX1; the conductor and
interval restrictions in QR25 remain in the HC5 application.\\
Actual frame and field, and admission&AH18 and HX2, with the exact
\texttt{HostAx} equality.\\
Pair rows, large $q,d$ and modulus bound&HC7: the given
\texttt{PairGood} and \texttt{PairWin}, with $Q\le C_0n$.\\
Five compatible auxiliary primes&HC5, applied with the actual $N=Qqd$;
the small primes lie below $e^{x^\varepsilon}$ and the norm primes
above it.\\
Second set and its two cardinality bounds&HC6 at the actual $L_w$;
one may use $C_{T_2}=64$.\\
Biased top set, shares and bad-prime bound&HC1--HC2, with the original
Frobenius steps and a record update changing only $TT$.\\
Real-character boxes and thin-character estimate&QP23--QP24 and
HC8--HC9 on that same host; $Q_{\mathrm{odd}}\mid M$ follows from
$Q\mid\operatorname{lcm}(2^{e+2},M)$.\\
Logarithmic derivative&HC8 gives $|dlogH|\le4\log x$; enlarge the
threshold until $4\log x\le x^\varepsilon$.\\
\end{tabular}
\end{center}

Choose $C_8=\max(C_0,C_1)$, $C_5=4$, and the supplied positive
$c_6,C_b$.  Take the maximum of the thresholds for these finitely
many results.  This is before every later $s,R,H_0,M,Y$.
The exact equality $H.\mathrm{toHostAx}=H_0$ preserves $Q,q,d,E$
and all the character data.  Installing HC2's $TT$ changes none of
them.  The norm identities give the required integral elements of
$\mathcal O_E\cap K$, congruent to $1$ modulo $Qqd$.
The congruences $p,p'\equiv1\pmod{\operatorname{lcm}(15,Qqd)}$
and $q_0\equiv1\pmod{Qqd}$ give all the splitting clauses by the
actual Artin map.  The level set is exactly
$\{g_0,g_1,p,p',q_0\}$.

The table supplies every conjunct of \texttt{HostGood}: its size and
arithmetic records, both second-set inequalities, \texttt{QDq},
\texttt{BoxW}, the logarithmic derivative, \texttt{ThinOK},
\texttt{BadOK}, disjoint $BB$ sets, integral norm elements, the two
congruence primes, the extra prime and the precise level set.
It also supplies the separate shares conclusion.  This closes the
literal \texttt{hostOfFieldOf} implication with no remaining analytic
input beyond the matched classical leaves used by QP and HX1.
The supplied \texttt{PairGood}, field, and rank inequalities remain
hypotheses of that implication.

\paragraph{HX4: install the full target request on the cohomological window.}
Assume $\mathrm{VWinAt}(\psi,401)$.  AH18 chooses integers
$c\ge1,e_0$ before the dyadic page and pair.  QP22 and QP25, applied
to that $c$ and the given window hypothesis, choose $J\ge1$ and
arbitrarily late intervals $[Z,Z^{3/2}]$ on which the literal
\texttt{WinDatum} exists.  Its construction uses a single $Y,q,d$
through each such interval, although the dyadic page $e=e(x)$ may
vary.  It has
\[
 \begin{gathered}
 x^{1/J}\le Y\le x,\qquad \mathrm{PairGood}(M,\varepsilon,Y,x,q,d),\\
 x^{2+8\varepsilon}\le c2^e\le2x^{2+8\varepsilon},\\
 \mathrm{VVle}\bigl(E,\lfloor x^{400}\rfloor,
                  \sqrt{\exp(x^\varepsilon)}\bigr)
 \quad\hbox{for every }\mathrm{IsChiField}(\psi,e,q,d,E).
 \end{gathered}
 \tag{HX.8}
\]
For precision, the passage from (HX.4) to (HX.8) uses
$\kappa=c^{1/(2+8\varepsilon)}$, $Z=\kappa W$ and
$x=\kappa x'$.  Taking
$W\ge\max(\kappa^3,\kappa^{400})$ gives
\[
 [Z,Z^{3/2}]\subseteq[\kappa W,W^2],\qquad
 W\le x'\le W^2,\qquad x^{400}\le(x')^{401}.
 \tag{HX.9}
\]
Also $x'\le x$, so (HX.2) preserves the bound
$\sqrt{\exp(x^\varepsilon)}$.  The exponent $400$ in (HX.8) is
therefore intentional; the later proof does not silently restore
$401$ after this rescaling.

Put
\[
 C_0=J+o+8M,\qquad K=5+4M,\qquad A=K^{5/\varepsilon}.
 \tag{HX.10}
\]
All are fixed before the scale.  Enlarge $x_0$ past $A^J$, the
rank and host thresholds, and $c2^{e_0}$.  Then
$Y\ge x^{1/J}\ge A$ gives $Y^{\varepsilon/5}\ge5+4M$.
The lower bound in (HX.8) forces $e\ge e_0$: otherwise
$c2^e\le c2^{e_0}<x\le x^{2+8\varepsilon}$.

The \texttt{PairGood} congruences give the exact \texttt{HostHyp},
$q\ne3$, coprimality with $M$, and the splitting conditions for
every prime at most $\lfloor Y^{\varepsilon/5}\rfloor$.
AH18 therefore supplies an actual frame $\phi,H_0$ with
\[
 \begin{gathered}
 |\phi|=c2^e,\quad\phi.o=o,\quad H_0.q=q,\quad H_0.d=d,\\
 M\mid Q\mid\operatorname{lcm}(2^{e+2},M),\qquad
 Q\le4M|\phi|,
 \end{gathered}
 \tag{HX.11}
\]
and all of its original \texttt{FieldRows}, kept-character,
exceptional-character, kernel and twist conclusions.  HX2 gives
the ranks over this very frame.  The rank window and (HX.11)
give $Q\le C_0R.n$ (the source uses the available factor $8M$),
while $J,o,M\le C_0$.  HX3 applies to these data, returning
$H.\mathrm{toHostAx}=H_0$ and the required $\mathrm{HostGood}$
and shares statements, with constants independent of this scale.

It remains to verify the two conjuncts of the target request;
neither is a new existence theorem.  For \texttt{EAtH}, its first
three fields are the equality $\phi.o=o$, kept-character statement
and exceptional-character statement from AH18.  Its shares field
is HX3.  For the last two fields we give the exact
\texttt{boundedTw} argument, extending GH4's finite-character
identification.  Let $\theta$ be in the source's bounded set
$(X^+\setminus\mathrm{Kept})\cup X^-$, and let $\chi$ be any
of its primitive Dirichlet representatives.  Both branches are
trivial on the first frame factor, so AH18's \texttt{TwistForm}
gives a power $\xi$ of $\psi$ and $i,j\in\{0,1\}$ with
\[
 H.\mathrm{dirAt}(\theta,a)=\xi(a)(a/q)^i(a/d)^j
       \quad((a,Qqd)=1),\qquad
 \theta(\iota_d)=(-1)^j.
 \tag{HX.12}
\]
The pair $(i,j)$ cannot be $(0,0)$.  In the $X^+$ branch, that
would make the primitive representative of $\xi$ one of the
kept characters and identify it with $\theta$, using $M\mid Q$;
this contradicts the definition of that branch.  In the $X^-$
branch, the displayed sign would be $+1$ instead of $-1$.

Every prime dividing $Q$ divides $2M$, by (HX.11).
Consequently an integer coprime to $2Mqd$ is coprime to $Qqd$.
The identity (HX.12), together with the defining identity of the
primitive representative $\chi$, therefore says exactly that
$\chi$ induces the displayed twist at the source level $2Mqd$.
This verifies the full \texttt{An.Induces} premise of
\texttt{PairGood.bounds}, not just an equality away from an
unspecified finite set of primes.  Apply that bound at $y=x$,
using $Y\le x$, to get
\[
 |G_\chi(x)|\le x^{1/2+\varepsilon/2}
                   \le x^{1/2+\varepsilon},\qquad
 |L'/L(1,\chi)|\le4\log x.
 \tag{HX.13}
\]
These are precisely the two remaining \texttt{EAtH} fields,
uniformly for every such primitive representative.

For cohomology, the same kernel identity identifies the actual
$H.E$ as $\mathrm{IsChiField}(\psi,e,q,d,H.E)$, so the last line of
(HX.8) applies to this field.  HC7 gives $R.t\le x^{400}$ after a
single further threshold; since $R.t$ is an integer,
$R.t\le\lfloor x^{400}\rfloor$.  Finally
$s.b=\exp(x^\varepsilon)$ by the scale definition.  Apply (HX.2)
to obtain exactly $\mathrm{VVle}(H.E,R.t,\sqrt{s.b})$.

All clauses of (HX.5) now hold.  Taking as scale set the union of
the \texttt{WinDatum} intervals proves the literal
$\mathrm{HostExW}(3/2,(\mathrm{GChi}\,\psi).\mathrm{req})$.
The constants may depend on the fixed character and $\varepsilon$;
they precede $x$, the page, the rank record, and the host.  The same
construction applies to $\psi=1$, with $M=o=c=1$; it does not
remove the cohomological premise (HX.4) from the RH specialization.

\paragraph{What is removed from the remaining dependency list.}
The actual target host is now an explicit consequence of the window
estimate, with the matching of all previous host inputs written out.
It is not an independent deep supplier.  In a connected proof graph
its incoming arithmetic edge is the complete quantified
\texttt{VWinAt} estimate, together with the classical leaves of
the established prime-pair/host constructions.  Constructing the
compatible theta/form/lattice family above that host is a separate
task, represented by \texttt{DesignExists}; HX4 does not prove it.
